\partieMat{1}{Plus, moins, autant}{Comparer des quantités jusqu’à six, et bien plus.}{\pic[scale=2.6] at (16.8,6.4) {panier};\pic[scale=1.5,mat/c=matRouge] at (16,7.4) {pomme};\pic[scale=1.5,mat/c=matVert] at (17.5,7.5) {pomme};\pic[scale=1.5,mat/c=matRouge] at (16.7,8.4) {pomme};\pic[scale=1.9,mat/c=matRouge] at (14,9.6) {pomme};\pic[scale=1.4,mat/c=matVert] at (18.8,9.8) {pomme};}{À partir de quatre ans, l’enfant compare des collections globalement, par correspondance terme à terme et en dénombrant chacune. Il comprend qu’une quantité ne change pas quand on déplace les objets.}{\ligneAct{1}{Cinq partout}{act-ms-1}\ligneAct{2}{Le plus de billes}{act-ms-2}\ligneAct{3}{Autant de jetons ?}{act-ms-3}\ligneAct{4}{Un verre par assiette}{act-ms-4}\ligneAct{5}{Voitures ou bateaux ?}{act-ms-5}}
\begin{activite}{1}{Cinq partout}{colorier}{Colorie le rond sous chaque boîte où il y a {\otBold\color{otNavy}cinq} objets.}{Reconnaître des collections de cinq objets de tailles et de dispositions différentes.}{act-ms-1}
\matFigurePage{ms-01}
\end{activite}
\begin{activite}{2}{Le plus de billes}{entourer}{Sans compter, entoure le bocal où il y a le {\otBold\color{otNavy}plus} de billes.}{Comparer globalement deux collections dont l’une a au moins deux fois plus d’objets que l’autre.}{act-ms-2}
\matFigurePage{ms-02}
\end{activite}
\begin{activite}{3}{Autant de jetons ?}{relier,dire}{Relie chaque jeton rouge à un jeton bleu. Dans chaque cadre, y a-t-il {\otBold\color{otNavy}autant} de jetons rouges que de bleus ?}{Comprendre qu’une quantité ne dépend pas de la place des objets : une rangée plus longue n’a pas forcément plus d’objets.}{act-ms-3}
\matFigurePage{ms-03}
\end{activite}
\begin{activite}{4}{Un verre par assiette}{colorier}{Colorie juste ce qu’il faut de verres pour mettre {\otBold\color{otNavy}un verre} à côté de chaque assiette.}{Dénombrer une collection pour en constituer une autre de même quantité, sans pouvoir les rapprocher.}{act-ms-4}
\matFigurePage{ms-04}
\end{activite}
\begin{activite}{5}{Voitures ou bateaux ?}{entourer,dire}{Compte les voitures, puis les bateaux. Entoure le groupe où il y en a le {\otBold\color{otNavy}plus}.}{Comparer deux collections en dénombrant chacune d’elles.}{act-ms-5}
\matFigurePage{ms-05}
\end{activite}
\partieMat{2}{Compter jusqu’à six}{Dénombrer, constituer, associer quantité et chiffre.}{\pic[scale=1.7,mat/n=3] at (14.6,9) {de};\pic[scale=1.7,mat/n=2] at (17.2,7) {de};\pic[scale=2.1,mat/n=3] at (18.6,9.6) {main};\pic[scale=1.3,mat/n=1] at (13.2,6.6) {de};}{L’enfant dénombre des collections jusqu’à six en pointant chaque objet une seule fois, utilise ses doigts et les constellations du dé, et associe chaque quantité à son écriture chiffrée.}{\ligneAct{6}{Le dé et les boîtes}{act-ms-6}\ligneAct{7}{Doigts et chiffres}{act-ms-7}\ligneAct{8}{Combien de points ?}{act-ms-8}\ligneAct{9}{Autant que le dé}{act-ms-9}\ligneAct{10}{Un ballon de plus}{act-ms-10}\ligneAct{11}{Autant que le chiffre}{act-ms-11}\ligneAct{12}{Le papillon butine}{act-ms-12}}
\begin{activite}{6}{Le dé et les boîtes}{relier}{Compte les objets de chaque boîte. Relie la boîte au dé qui a {\otBold\color{otNavy}autant} de points.}{Dénombrer une collection jusqu’à six et l’associer à la constellation du dé.}{act-ms-6}
\matFigurePage{ms-06}
\end{activite}
\begin{activite}{7}{Doigts et chiffres}{relier}{Compte les doigts levés. Relie chaque dessin au {\otBold\color{otNavy}chiffre} qui convient.}{Associer une quantité montrée avec les doigts et son écriture chiffrée.}{act-ms-7}
\matFigurePage{ms-07}
\end{activite}
\begin{activite}{8}{Combien de points ?}{entourer}{Compte les points de chaque carte. Entoure le {\otBold\color{otNavy}bon chiffre}.}{Dénombrer une collection mal rangée en pointant chaque élément une seule fois.}{act-ms-8}
\matFigurePage{ms-08}
\end{activite}
\begin{activite}{9}{Autant que le dé}{dessiner}{Dessine dans chaque boîte {\otBold\color{otNavy}autant} de pommes qu’il y a de points sur le dé.}{Constituer une collection de même quantité qu’une constellation donnée.}{act-ms-9}
\matFigurePage{ms-09}
\end{activite}
\begin{activite}{10}{Un ballon de plus}{dessiner,dire}{Dessine {\otBold\color{otNavy}un ballon de plus} dans la case vide. Dis combien il y en avait et combien il y en a maintenant.}{Réaliser une collection qui a un objet de plus et nommer le nombre avant et après l’ajout.}{act-ms-10}
\matFigurePage{ms-10}
\end{activite}
\begin{activite}{11}{Autant que le chiffre}{colorier}{Colorie {\otBold\color{otNavy}autant} de cases que le chiffre écrit au début de la ligne.}{Passer de l’écriture chiffrée d’un nombre à la quantité.}{act-ms-11}
\matFigurePage{ms-11}
\end{activite}
\begin{activite}{12}{Le papillon butine}{tracer}{Trace le chemin du papillon : il doit passer sur {\otBold\color{otNavy}chaque fleur une seule fois}, sans en oublier.}{Parcourir une collection en passant une et une seule fois par chaque élément.}{act-ms-12}
\matFigurePage{ms-12}
\end{activite}
\partieMat{3}{Composer et écrire}{Les nombres jusqu’à six ; la comptine jusqu’à douze.}{\pic[scale=1.7,mat/n=5] at (14,9.2) {de};\pic[scale=2.3,mat/n=4] at (17.8,9.2) {main};\pic[scale=1.4,mat/n=6] at (16,6.6) {de};}{Les compositions et décompositions (cinq, c’est deux et trois) permettent de dénombrer plus vite. L’enfant s’initie à l’écriture des chiffres de un à six et récite la comptine jusqu’à douze.}{\ligneAct{13}{Les doigts des deux mains}{act-ms-13}\ligneAct{14}{Cinq perles, deux couleurs}{act-ms-14}\ligneAct{15}{J’écris 1, 2, 3}{act-ms-15}\ligneAct{16}{J’écris 4, 5, 6}{act-ms-16}\ligneAct{17}{J’écris combien}{act-ms-17}\ligneAct{18}{Les pierres de la rivière}{act-ms-18}}
\begin{activite}{13}{Les doigts des deux mains}{entourer,dire}{Combien de doigts sont levés en tout ? Entoure le {\otBold\color{otNavy}bon chiffre}.}{Composer des nombres jusqu’à six avec les doigts des deux mains.}{act-ms-13}
\matFigurePage{ms-13}
\end{activite}
\begin{activite}{14}{Cinq perles, deux couleurs}{colorier,dire}{Colorie chaque collier en rouge et en bleu. Fais {\otBold\color{otNavy}chaque fois} un collier différent.}{Décomposer cinq de plusieurs façons : un et quatre, deux et trois…}{act-ms-14}
\matFigurePage{ms-14}
\end{activite}
\begin{activite}{15}{J’écris 1, 2, 3}{tracer}{Repasse sur les chiffres gris avec ton doigt, puis avec un crayon. Écris le chiffre {\otBold\color{otNavy}tout seul} dans la dernière case.}{Écrire en chiffres les nombres de un à six.}{act-ms-15}
\matFigurePage{ms-15}
\end{activite}
\begin{activite}{16}{J’écris 4, 5, 6}{tracer}{Repasse sur les chiffres gris avec ton doigt, puis avec un crayon. Écris le chiffre {\otBold\color{otNavy}tout seul} dans la dernière case.}{Écrire en chiffres les nombres de un à six.}{act-ms-16}
\matFigurePage{ms-16}
\end{activite}
\begin{activite}{17}{J’écris combien}{dessiner}{Compte les animaux ou les étoiles. {\otBold\color{otNavy}Écris le chiffre} dans la case.}{Représenter une quantité par son écriture chiffrée.}{act-ms-17}
\matFigurePage{ms-17}
\end{activite}
\begin{activite}{18}{Les pierres de la rivière}{dire,entourer}{Fais sauter le lapin de pierre en pierre avec ton doigt en disant : un, deux, trois… Entoure la dernière pierre : c’est la {\otBold\color{otNavy}douzième}.}{Réciter la comptine numérique de un à douze, de façon ordonnée, un mot par pierre.}{act-ms-18}
\matFigurePage{ms-18}
\end{activite}
\partieMat{4}{Rangs et cases}{Le premier, le deuxième… la bande numérique et le jeu de l’oie.}{\pic[scale=1.8,mat/c=matRouge] at (14,8) {pion};\pic[scale=1.8,mat/c=matBleu] at (15.8,7.2) {pion};\pic[scale=1.6,mat/n=4] at (18.2,9.2) {de};}{Le nombre sert aussi à repérer une position : un rang dans une file, une case sur une piste. L’enfant compte à partir d’un point de départ, avance ou recule d’une ou deux cases, et découvre le début de la bande numérique.}{\ligneAct{19}{En partant de la mare}{act-ms-19}\ligneAct{20}{Le train des rangs}{act-ms-20}\ligneAct{21}{La perle rouge}{act-ms-21}\ligneAct{22}{La bande des nombres}{act-ms-22}\ligneAct{23}{Chaque carte à sa place}{act-ms-23}\ligneAct{24}{Le jeu de l’oie}{act-ms-24}\ligneAct{25}{On recule !}{act-ms-25}}
\begin{activite}{19}{En partant de la mare}{entourer,barrer}{En partant de la mare, entoure le {\otBold\color{otNavy}deuxième} animal. Fais une croix sur le {\otBold\color{otNavy}cinquième}.}{Repérer le rang d’un objet dans une file en comptant, à partir d’un point de départ donné.}{act-ms-19}
\matFigurePage{ms-19}
\end{activite}
\begin{activite}{20}{Le train des rangs}{colorier}{Colorie le {\otBold\color{otNavy}premier} wagon en rouge, le {\otBold\color{otNavy}troisième} en bleu et le {\otBold\color{otNavy}sixième} en vert. On compte à partir de la locomotive.}{Repérer un rang jusqu’au sixième dans une suite ordonnée.}{act-ms-20}
\matFigurePage{ms-20}
\end{activite}
\begin{activite}{21}{La perle rouge}{entourer}{En partant du nœud, entoure le fil où la perle rouge est la {\otBold\color{otNavy}troisième}.}{Repérer le rang d’un élément en comptant à partir d’un point de départ.}{act-ms-21}
\matFigurePage{ms-21}
\end{activite}
\begin{activite}{22}{La bande des nombres}{dessiner}{Écris les nombres qui {\otBold\color{otNavy}manquent} dans chaque bande.}{Compléter une bande numérique lacunaire jusqu’à six.}{act-ms-22}
\matFigurePage{ms-22}
\end{activite}
\begin{activite}{23}{Chaque carte à sa place}{relier}{Relie chaque carte à la case de la bande qui a le {\otBold\color{otNavy}même nombre}.}{Placer des représentations des nombres (dé, doigts, points) sur la bande numérique.}{act-ms-23}
\matFigurePage{ms-23}
\end{activite}
\begin{activite}{24}{Le jeu de l’oie}{colorier}{Le pion avance d’autant de cases que de points sur le dé. Colorie la case où il {\otBold\color{otNavy}arrive}.}{Trouver la position d’arrivée après un avancement d’une ou deux cases.}{act-ms-24}
\matFigurePage{ms-24}
\end{activite}
\begin{activite}{25}{On recule !}{colorier}{Attention : la flèche rouge veut dire {\otBold\color{otNavy}reculer}. Colorie la case où le pion arrive.}{Trouver la position d’arrivée après un recul d’une ou deux cases.}{act-ms-25}
\matFigurePage{ms-25}
\end{activite}
\partieMat{5}{Petits problèmes}{Ajouter, enlever, cacher, regrouper, partager.}{\pic[scale=2.6,mat/c=matBrun] at (15,8) {ours};\pic[scale=2.4,mat/c=matOrange] at (18.2,6.8) {valise};\pic[scale=1.5,mat/c=matJaune] at (18.6,9.8) {ours};}{Les problèmes se jouent d’abord avec des objets, puis se résolvent sur la fiche. L’enfant peut compter sur ses doigts, surcompter ou utiliser une décomposition (« six, c’est quatre et deux »).}{\ligneAct{26}{Les fraises de Lilou}{act-ms-26}\ligneAct{27}{Le chapeau du magicien}{act-ms-27}\ligneAct{28}{Les billes dans la boîte}{act-ms-28}\ligneAct{29}{Sous la main}{act-ms-29}\ligneAct{30}{Les boîtes de crayons}{act-ms-30}\ligneAct{31}{Le partage des gâteaux}{act-ms-31}\ligneAct{32}{Une part pour chaque lapin}{act-ms-32}}
\begin{activite}{26}{Les fraises de Lilou}{barrer,entourer}{Lilou avait cinq fraises. Elle en a mangé deux. {\otBold\color{otNavy}Barre} les fraises mangées, puis entoure le dé qui montre combien il en reste.}{Résoudre un problème de retrait en représentant l’action.}{act-ms-26}
\matFigurePage{ms-26}
\end{activite}
\begin{activite}{27}{Le chapeau du magicien}{dessiner}{Il y a six lapins en tout. Certains sont cachés sous le chapeau. {\otBold\color{otNavy}Dessine} dans la boîte autant de ronds que de lapins cachés.}{Trouver une partie quand on connaît le tout et l’autre partie.}{act-ms-27}
\matFigurePage{ms-27}
\end{activite}
\begin{activite}{28}{Les billes dans la boîte}{entourer}{Il y a trois billes dans la boîte. On en met encore deux. Combien y a-t-il de billes dans la boîte ? {\otBold\color{otNavy}Entoure} le bon chiffre.}{Résoudre un problème d’ajout : trouver le tout.}{act-ms-28}
\matFigurePage{ms-28}
\end{activite}
\begin{activite}{29}{Sous la main}{dessiner}{Le dé dit combien de jetons il y a en tout. Combien sont cachés sous la main ? {\otBold\color{otNavy}Dessine-les} dans la boîte.}{Rechercher une partie dans un problème de parties-tout.}{act-ms-29}
\matFigurePage{ms-29}
\end{activite}
\begin{activite}{30}{Les boîtes de crayons}{entourer}{Il y a deux crayons dans chaque boîte. Combien de crayons y a-t-il {\otBold\color{otNavy}en tout} ? Entoure le bon dé.}{Rechercher le tout dans un problème de groupements.}{act-ms-30}
\matFigurePage{ms-30}
\end{activite}
\begin{activite}{31}{Le partage des gâteaux}{dessiner}{Partage les six gâteaux entre les deux ours : chacun doit en avoir {\otBold\color{otNavy}autant}. Dessine les gâteaux dans les assiettes.}{Trouver la valeur d’une part dans un partage équitable, en distribuant un à un.}{act-ms-31}
\matFigurePage{ms-31}
\end{activite}
\begin{activite}{32}{Une part pour chaque lapin}{dessiner}{Partage les six carottes entre les trois lapins : chacun doit en avoir {\otBold\color{otNavy}autant}. Dessine-les dans les cases.}{Partager équitablement une collection en distribuant un à un.}{act-ms-32}
\matFigurePage{ms-32}
\end{activite}
\partieMat{6}{Formes et solides}{Triangles, carrés, disques ; empreintes et assemblages.}{\draw[mat contour,fill=matRouge] (15,8) circle (1.2);\draw[mat contour,fill=matJaune] (17.565,8.204) -- (19.196,8.965) -- (18.435,10.596) -- (16.804,9.835) -- cycle;\draw[mat contour,fill=otGreen] (16.815,4.533) -- (19.11,5.138) -- (17.392,6.645) -- cycle;}{L’enfant reconnaît et classe les formes planes dans toutes les orientations. Les solides (cube, boule, pyramide, cylindre) se manipulent : la fiche travaille leurs empreintes. Aucun solide n’est dessiné en perspective.}{\ligneAct{33}{Triangles, carrés, disques}{act-ms-33}\ligneAct{34}{L’intrus}{act-ms-34}\ligneAct{35}{Les empreintes du cube}{act-ms-35}\ligneAct{36}{Les empreintes de la pyramide}{act-ms-36}\ligneAct{37}{La fusée}{act-ms-37}\ligneAct{38}{Le carré en cinq pièces}{act-ms-38}}
\begin{activite}{33}{Triangles, carrés, disques}{colorier}{Colorie les triangles en {\otBold\color{otNavy}bleu}, les carrés en {\otBold\color{otNavy}rouge} et les disques en {\otBold\color{otNavy}jaune}.}{Reconnaître et classer triangles, carrés et disques, quelles que soient leur taille et leur orientation.}{act-ms-33}
\matFigurePage{ms-33}
\end{activite}
\begin{activite}{34}{L’intrus}{barrer}{Dans chaque ligne, une forme n’est pas comme les autres. {\otBold\color{otNavy}Barre-la}.}{Reconnaître une forme plane malgré sa couleur, sa taille et son orientation.}{act-ms-34}
\matFigurePage{ms-34}
\end{activite}
\begin{activite}{35}{Les empreintes du cube}{manipuler,entourer}{Avec un vrai cube et de la pâte à modeler, fais des empreintes. {\otBold\color{otNavy}Entoure} les empreintes que le cube peut laisser.}{Produire les empreintes d’un solide et les reconnaître.}{act-ms-35}
\matFigurePage{ms-35}
\end{activite}
\begin{activite}{36}{Les empreintes de la pyramide}{manipuler,entourer}{Avec une vraie pyramide, fais des empreintes. {\otBold\color{otNavy}Entoure} les empreintes que la pyramide peut laisser.}{Produire les empreintes d’un solide et retrouver le solide à partir d’une empreinte.}{act-ms-36}
\matFigurePage{ms-36}
\end{activite}
\begin{activite}{37}{La fusée}{colorier}{Colorie la grande fusée comme le {\otBold\color{otNavy}modèle}.}{Reproduire un assemblage de cinq formes à partir d’un modèle.}{act-ms-37}
\matFigurePage{ms-37}
\end{activite}
\begin{activite}{38}{Le carré en cinq pièces}{colorier}{Colorie le grand carré comme le {\otBold\color{otNavy}modèle}. Regarde bien la place de chaque pièce.}{Reproduire un pavage de cinq pièces à partir d’un modèle.}{act-ms-38}
\matFigurePage{ms-38}
\end{activite}
\partieMat{7}{Longueurs et masses}{Plus long, plus court ; plus lourd, plus léger.}{\pic[scale=2.4,mat/n=-1] at (16,8.2) {balance};\pic[scale=1.2] at (13.5,9.9) {pasteque};\pic[scale=1.3] at (18.6,10.5) {plume};}{L’enfant classe et range des objets selon leur longueur, en les mettant à la même origine quand c’est nécessaire. La masse est introduite à partir de quatre ans : on soupèse, puis on utilise une balance de Roberval.}{\ligneAct{39}{Courts, moyens, longs}{act-ms-39}\ligneAct{40}{Du plus court au plus long}{act-ms-40}\ligneAct{41}{De même longueur}{act-ms-41}\ligneAct{42}{Le plus lourd}{act-ms-42}\ligneAct{43}{La balance}{act-ms-43}}
\begin{activite}{39}{Courts, moyens, longs}{colorier}{Colorie les rubans courts en {\otBold\color{otNavy}jaune}, les moyens en {\otBold\color{otNavy}vert} et les longs en {\otBold\color{otNavy}rouge}.}{Classer des objets selon leur longueur en trois groupes, quand les longueurs sont très différentes.}{act-ms-39}
\matFigurePage{ms-39}
\end{activite}
\begin{activite}{40}{Du plus court au plus long}{dessiner}{Dessine {\otBold\color{otNavy}un} point à côté du crayon le plus court, {\otBold\color{otNavy}deux} points à côté du suivant… et {\otBold\color{otNavy}quatre} points à côté du plus long.}{Ranger des objets rectilignes du plus court au plus long et verbaliser le résultat.}{act-ms-40}
\matFigurePage{ms-40}
\end{activite}
\begin{activite}{41}{De même longueur}{entourer}{Dans chaque cadre, {\otBold\color{otNavy}entoure} les deux rubans qui ont la même longueur.}{Comparer des longueurs et reconnaître deux objets de même longueur.}{act-ms-41}
\matFigurePage{ms-41}
\end{activite}
\begin{activite}{42}{Le plus lourd}{manipuler,entourer}{Dans chaque cadre, {\otBold\color{otNavy}entoure} l’objet le plus lourd. Avant, soupèse de vrais objets.}{Comparer les masses de deux objets très différents ; comprendre qu’un gros objet n’est pas forcément lourd.}{act-ms-42}
\matFigurePage{ms-42}
\end{activite}
\begin{activite}{43}{La balance}{entourer}{Le plateau le plus lourd descend. Sur chaque balance, {\otBold\color{otNavy}entoure} l’objet le plus lourd.}{Utiliser une balance de Roberval pour comparer deux masses.}{act-ms-43}
\matFigurePage{ms-43}
\end{activite}
\partieMat{8}{Les motifs}{Continuer, trouver l’intrus, mémoriser.}{\pic[scale=1.4,mat/c=matRouge] at (12,7.8) {perle};\pic[scale=1.4,mat/c=matBleu] at (13.1,8.3) {perle};\pic[scale=1.4,mat/c=matRouge] at (14.2,7.8) {perle};\pic[scale=1.4,mat/c=matBleu] at (15.3,8.3) {perle};\pic[scale=1.4,mat/c=matRouge] at (16.4,7.8) {perle};\pic[scale=1.4,mat/c=matBleu] at (17.5,8.3) {perle};\pic[scale=1.4,mat/c=matRouge] at (18.6,7.8) {perle};}{Un motif répétitif suit une règle : l’enfant la repère, la dit avec ses mots, trouve l’intrus et prolonge le motif. Toute proposition cohérente est acceptée si l’enfant explique sa règle.}{\ligneAct{44}{Continue la rangée}{act-ms-44}\ligneAct{45}{L’intrus de la rangée}{act-ms-45}\ligneAct{46}{Ça se répète ?}{act-ms-46}\ligneAct{47}{La tour de mémoire}{act-ms-47}\ligneAct{48}{Sous la carte}{act-ms-48}}
\begin{activite}{44}{Continue la rangée}{dessiner,dire}{Regarde comment chaque rangée est faite. {\otBold\color{otNavy}Dessine} les trois formes qui viennent après.}{Prolonger un motif répétitif et dire la règle utilisée.}{act-ms-44}
\matFigurePage{ms-44}
\end{activite}
\begin{activite}{45}{L’intrus de la rangée}{barrer}{Dans chaque rangée, une forme n’est pas à sa place. {\otBold\color{otNavy}Barre-la}.}{Reconnaître un motif répétitif à ses régularités et trouver l’élément qui ne respecte pas la règle.}{act-ms-45}
\matFigurePage{ms-45}
\end{activite}
\begin{activite}{46}{Ça se répète ?}{entourer}{Entoure les rangées où les formes {\otBold\color{otNavy}se répètent} toujours de la même façon.}{Identifier parmi plusieurs rangées celles qui contiennent un motif répétitif.}{act-ms-46}
\matFigurePage{ms-46}
\end{activite}
\begin{activite}{47}{La tour de mémoire}{colorier}{Regarde bien la tour du modèle. {\otBold\color{otNavy}Cache-la} avec ta main, puis colorie l’autre tour de mémoire.}{Mémoriser un motif répétitif simple pour le reproduire.}{act-ms-47}
\matFigurePage{ms-47}
\end{activite}
\begin{activite}{48}{Sous la carte}{entourer}{Quelle forme est cachée sous la carte bleue ? {\otBold\color{otNavy}Entoure-la} dans le cadre du dessous.}{Compléter un motif répétitif en trouvant l’élément manquant.}{act-ms-48}
\matFigurePage{ms-48}
\end{activite}
\partieMat{9}{Se repérer}{Entre, au milieu, à l’intérieur ; suivre un chemin.}{\pic[scale=2.2] at (14.6,7) {souris};\pic[scale=2] at (18.4,6.6) {fromage};\pic[scale=2] at (16.8,9.6) {arbre};}{L’enfant situe les objets les uns par rapport aux autres et utilise un vocabulaire plus riche : entre, au milieu, à l’intérieur, à l’extérieur. Il suit un déplacement sur un dessin et le raconte.}{\ligneAct{49}{Entre les deux}{act-ms-49}\ligneAct{50}{En haut, au milieu, en bas}{act-ms-50}\ligneAct{51}{Dans l’enclos}{act-ms-51}\ligneAct{52}{Le labyrinthe}{act-ms-52}\ligneAct{53}{Le chemin du chat}{act-ms-53}}
\begin{activite}{49}{Entre les deux}{entourer}{En haut, entoure le jouet qui est {\otBold\color{otNavy}entre} le chat et le lapin. En bas, entoure celui qui est entre le bateau et le ballon.}{Situer des objets les uns par rapport aux autres : « entre ».}{act-ms-49}
\matFigurePage{ms-49}
\end{activite}
\begin{activite}{50}{En haut, au milieu, en bas}{entourer,barrer,colorier}{{\otBold\color{otNavy}Entoure} le jouet tout en haut, {\otBold\color{otNavy}barre} celui tout en bas et colorie celui du milieu.}{Utiliser les mots « en haut », « en bas », « au milieu » pour situer des objets.}{act-ms-50}
\matFigurePage{ms-50}
\end{activite}
\begin{activite}{51}{Dans l’enclos}{colorier}{Colorie les animaux qui sont {\otBold\color{otNavy}à l’intérieur} de l’enclos.}{Comprendre et utiliser « à l’intérieur » et « à l’extérieur ».}{act-ms-51}
\matFigurePage{ms-51}
\end{activite}
\begin{activite}{52}{Le labyrinthe}{tracer}{Trace le chemin de la souris jusqu’au fromage, {\otBold\color{otNavy}sans traverser} les murs.}{Suivre un déplacement sur une représentation et le tracer.}{act-ms-52}
\matFigurePage{ms-52}
\end{activite}
\begin{activite}{53}{Le chemin du chat}{dire,tracer}{Suis la route avec ton doigt, du chat jusqu’à sa maison. {\otBold\color{otNavy}Raconte} le trajet, puis trace-le.}{Raconter un déplacement à partir d’une représentation, avec des repères et les mots de l’espace.}{act-ms-53}
\matFigurePage{ms-53}
\end{activite}
